\chapter{Combining the bases and applying the classical test}
\label{chap:qpbt-combining}
% Paper subsection label sec:combining kept for cross-reference fidelity; the
% companion label sec:apply-ldt is attached to the corresponding section below.
% Source: references/qpbt-paper/14_analysis_of_the_pauli_basis_test.tex,
% lines 680-1414 (subsections "Combining the X and Z measurements" and
% "Applying the classical low-degree test").
\label{sec:combining}

Throughout this chapter we work in the setting of
Chapter~\ref{chap:qpbt-observables}: a projective strategy for the Pauli basis
test game succeeds with probability at least $1-\eps$, and all operators act
relative to the expanded state $\ket{\hat\psi}$ on registers
$\mathsf{A}\,\mathsf{A}'\,\mathsf{A}''\,\mathsf{B}\,\mathsf{B}'\,\mathsf{B}''$
(Definition~\ref{def:expanded-state}), where $n$ denotes the number of ancilla
qudit registers, each isomorphic to $\C^q$, appended on each side. We write
$q = 2^t$ (Definition~\ref{def:admissible-size}) and $\tr:\F_q\to\F_2$ for the
subfield trace. The points measurements
$\hat{M}^{(\mathrm{Point},W),u}_a$ for $W \in \{X,Z\}$
(Definition~\ref{def:expanded-point-measurement}), their binary refinements
\[
	\hat{M}^{(\mathrm{Point},W),u,r}_b
	= \sum_{a \,:\, \tr(ar) = b} \hat{M}^{(\mathrm{Point},W),u}_a\,,
	\qquad r \in \F_q,\ b \in \F_2
\]
(Definition~\ref{def:expanded-point-trace-projection}), and the lines
measurements $\hat{M}^{(\mathrm{Line},W),\ell}_{f}$
(Definition~\ref{def:expanded-line-measurement}) are those constructed in
Chapter~\ref{chap:qpbt-observables}, and $\delta_{\mathrm{acom}}(\eps)$ and
$\delta_{\mathrm{Line}}(\eps)$ denote the error functions of
Lemma~\ref{lem:qld-comm-cons} and Lemma~\ref{lem:qld-comm-line-cons},
respectively.

A subscript such as $(\,\cdot\,)_{\mathsf{A}\mathsf{A}'}$ indicates the pair of
registers on which an operator acts; the two bipartitions in use are
$\mathsf{A}\mathsf{A}'$ versus $\mathsf{B}\mathsf{A}''$ and
$\mathsf{B}\mathsf{B}'$ versus $\mathsf{A}\mathsf{B}''$. The state-dependent
distances $\approx_\delta$ (Definition~\ref{def:povm-distance}) and
$\simeq_\delta$ (Definition~\ref{def:consistency}) are taken with respect to
$\ket{\hat\psi}$ and averaged over the stated question distribution, and
``all symmetric equivalents'' of an approximation is meant in the sense of
Remark~\ref{def:symmetric-equivalents}. Between scalars, $\alpha
\approx_\delta \beta$ means $|\alpha - \beta| \le O(\delta)$. For sums of
measurement operators over outcomes satisfying a predicate we use the bracket
notation of Definition~\ref{def:bracket}; for instance
$T^{\ell_X,\ell_Z}_{[\mathrm{eval}_x(\cdot)=a,\ \mathrm{eval}_z(\cdot)=b]}
= \sum_{f_X,f_Z \,:\, f_X(x)=a,\ f_Z(z)=b} T^{\ell_X,\ell_Z}_{f_X,f_Z}$.
For a line $\ell$ (Definition~\ref{def:line}) and $k \in \N$ we write
$\deg_k(\ell)$ for the set of polynomial functions on $\ell$ of degree at most
$k$ in the line parameter, and
$\mathrm{ideg}_{d,m}(\F_q) = \polyfunc{m}{q}{d}$
(Definition~\ref{def:polyfunc}) for the $m$-variate polynomial representatives
over $\F_q$ of individual degree at most $d$, identified with the induced
polynomial functions via evaluation (Definition~\ref{def:polynomials-degree};
the identification is bijective for $d \le q-1$, which holds throughout this
chapter's parameter regime); the restriction of an element of
$\mathrm{ideg}_{d,m}(\F_q)$ to a line has total degree at most $md$.

\section{Combined point measurements}

The measurements $\hat{M}^{(\mathrm{Point},X),x}$ and
$\hat{M}^{(\mathrm{Point},Z),z}$ commute only approximately. The first goal of
this chapter is to construct a single projective measurement that
simultaneously measures both, first at points and then along lines. The
construction of the combined point measurement proceeds in three steps: the
binary refinements $\hat{M}^{(\mathrm{Point},X),x,r}$ and
$\hat{M}^{(\mathrm{Point},Z),z,s}$ are combined into a sandwich POVM, the
resulting family is shown to satisfy an approximate linearity in the index
pair $(r,s)$, and the following result of Natarajan and
Vidick~\cite{Natarajan2017Linearity} converts approximate linearity into exact
linearity, from which a projective measurement with outcomes in $\F_q^2$ is
assembled by Fourier averaging.

\begin{theorem}[Quantum linearity test~\cite{Natarajan2017Linearity}]
	\label{thm:linearity}
	\uses{def:povm-distance}
	Let $0 \le \delta \le 1$, let $t$ be a positive integer, and let
	$\ket{\psi}$ be a state in a Hilbert space
	$\hilb_{\mathsf{A}} \ot \hilb_{\mathsf{B}}$. Let $\{\mathcal{O}^u\}$ be a
	set of observables indexed by $u \in \F_2^t$ acting on
	$\hilb_{\mathsf{A}}$, and suppose that on average over $u, u'$ chosen
	uniformly at random from $\F_2^t$ the approximate linearity
	\[
		(\mathcal{O}^u \, \mathcal{O}^{u'})_{\mathsf{A}}
		\approx_\delta (\mathcal{O}^{u+u'})_{\mathsf{A}}
	\]
	holds with respect to $\ket{\psi}$. Then there exist an extended state
	$\ket{\psi'}_{\mathsf{A}\mathsf{A}'\mathsf{B}}
	= \ket{\psi}_{\mathsf{A}\mathsf{B}} \ot \ket{0}_{\mathsf{A}'}$
	in an extended space
	$\hilb_{\mathsf{A}} \ot \hilb_{\mathsf{B}} \ot \hilb_{\mathsf{A}'}$ and
	observables $\{\mathcal{L}^u\}$ acting on
	$\hilb_{\mathsf{A}} \ot \hilb_{\mathsf{A}'}$ such that for all
	$u, u' \in \F_2^t$,
	\[
		\mathcal{L}^u \, \mathcal{L}^{u'} = \mathcal{L}^{u+u'}
		\qquad \text{and} \qquad
		(\mathcal{L}^u)_{\mathsf{A}\mathsf{A}'}
		\approx_\delta (\mathcal{O}^u)_{\mathsf{A}}\,,
	\]
	the latter with respect to $\ket{\psi'}$. Furthermore,
	$\mathcal{L}^u = \Id$ when $u = 0$.
\end{theorem}
\begin{proof}
	This is the quantum linearity test of Natarajan and
	Vidick~\cite{Natarajan2017Linearity}; no proof is given in
	\cite{Ji2020MIPStar}, which invokes the result directly. The statement is
	imported here as an external result; see
	Remark~\ref{rem:linearity-import}.
\end{proof}

\begin{remark}[Properties of the quantum linearity theorem used below]
	\label{rem:linearity-import}
	Theorem~\ref{thm:linearity} is the quantum linearity theorem of
	Natarajan and Vidick~\cite{Natarajan2017Linearity}, used here as an
	external result. Two features of the statement are load-bearing for the
	proof of Lemma~\ref{lem:qld-4-10}: the error is preserved exactly
	($\delta$ in the hypothesis equals $\delta$ in the conclusion), and the
	linearity of the $\{\mathcal{L}^u\}$ is exact and holds for all pairs
	$(u,u')$, not merely on average.
\end{remark}

\begin{lemma}[Combined point measurements]
	\label{lem:qld-4-10}
	\uses{def:expanded-state, def:expanded-point-measurement,
		def:povm-distance, def:symmetric-equivalents}
	There exists a function $\delta_Q(\eps) = \poly(\eps)$ such that the
	following holds. For every $x, z \in \F_q^m$ and
	$\hilb \in \{\hilb_{\mathsf{A}}, \hilb_{\mathsf{B}}\}$ there are
	projective measurements $\{\hat{Q}^{x,z}_{a,b}\}_{a,b \in \F_q}$ acting
	on $\hilb \ot (\C^q)^{\ot n}$ such that:
	\begin{enumerate}
		\item (\textbf{Self-consistency}) On average over uniformly random
		$x, z \in \F_q^m$,
		\begin{equation}
			\label{eq:qld-q-self-cons}
			(\hat{Q}^{x,z}_{a,b})_{\mathsf{A}\mathsf{A}'}
			\approx_{\delta_Q}
			(\hat{Q}^{x,z}_{a,b})_{\mathsf{B}\mathsf{A}''}\,.
		\end{equation}
		\item (\textbf{Consistency with $\hat{M}$}) On average over uniformly
		random $x, z \in \F_q^m$,
		\begin{gather}
			\label{eq:qld-q-cons-m-hat-xz}
			(\hat{Q}^{x,z}_{a,b})_{\mathsf{A}\mathsf{A}'}
			\approx_{\delta_Q}
			(\hat{M}^{(\mathrm{Point},X),x}_{a} \,
			\hat{M}^{(\mathrm{Point},Z),z}_{b})_{\mathsf{B}\mathsf{A}''}\,,
			\\
			\label{eq:qld-q-cons-m-hat-zx}
			(\hat{Q}^{x,z}_{a,b})_{\mathsf{A}\mathsf{A}'}
			\approx_{\delta_Q}
			(\hat{M}^{(\mathrm{Point},Z),z}_{b} \,
			\hat{M}^{(\mathrm{Point},X),x}_{a})_{\mathsf{B}\mathsf{A}''}\,.
		\end{gather}
	\end{enumerate}
	Furthermore, all symmetric equivalents of these approximations also hold.
	The operators $\hat{M}^{(\mathrm{Point},X),x}_{a}$ and
	$\hat{M}^{(\mathrm{Point},Z),z}_{b}$ do not commute in general, so the
	products in \eqref{eq:qld-q-cons-m-hat-xz} and
	\eqref{eq:qld-q-cons-m-hat-zx} differ and the two relations are separate
	conclusions.
\end{lemma}
\begin{proof}
	\uses{lem:qld-comm-cons, def:expanded-point-trace-projection,
		fact:add-a-proj, fact:agreement, lem:ortho, thm:linearity,
		lem:avg-closeness, lem:cancellation}
	Throughout the proof every intermediate error is a $\poly(\eps)$ function
	that may differ from one use of $\approx_\delta$ to the next. For
	$\omega = (x,z,r,s) \in \F_q^m \times \F_q^m \times \F_q \times \F_q$
	and $a, b \in \F_2$ define the sandwich POVM
	$\hat{R}^{\omega}_{a,b}
	= \hat{M}^{(\mathrm{Point},Z),z,s}_{b}
	\hat{M}^{(\mathrm{Point},X),x,r}_{a}
	\hat{M}^{(\mathrm{Point},Z),z,s}_{b}$
	from the binary refinements of
	Definition~\ref{def:expanded-point-trace-projection}; by the approximate
	commutation of Lemma~\ref{lem:qld-comm-cons}, projectivity, and
	Lemma~\ref{fact:add-a-proj}, on average over $\omega$ one has
	$\hat{R}^{\omega}_{a,b}
	\approx_{\sqrt{\delta_{\mathrm{acom}}}}
	\hat{M}^{(\mathrm{Point},Z),z,s}_{b}
	\hat{M}^{(\mathrm{Point},X),x,r}_{a}$,
	and a Cauchy--Schwarz chain with successive errors
	$\delta_{\mathrm{acom}}^{1/4}$, $\delta_{\mathrm{acom}}^{1/4}$,
	$\sqrt{\eps}$, $\eps$ establishes the self-consistency
	$(\hat{R}^{\omega}_{a,b})_{\mathsf{A}\mathsf{A}'} \simeq_\delta
	(\hat{R}^{\omega}_{a,b})_{\mathsf{B}\mathsf{A}''}$, hence the
	corresponding $\approx_\delta$ relation by Lemma~\ref{fact:agreement}.
	The orthonormalization lemma (Lemma~\ref{lem:ortho}) then produces
	projective measurements $\{\hat{L}^{\omega}_{a,b}\}$ with
	$\hat{L}^{\omega}_{a,b} \approx_\delta \hat{R}^{\omega}_{a,b}$ on
	average over $\omega$. Next, the refined point measurements satisfy the
	exact linearity
	$\sum_{b+b'=c} \hat{M}^{(\mathrm{Point},W),u,r}_{b}
	\hat{M}^{(\mathrm{Point},W),u,r'}_{b'}
	= \hat{M}^{(\mathrm{Point},W),u,r+r'}_{c}$
	(a direct consequence of projectivity of
	$\{\hat{M}^{(\mathrm{Point},W),u}_a\}$), which transfers, on average over
	$x, z, r, r', s, s'$, to an approximate linearity of the
	$\{\hat{R}^{\omega}_{a,b}\}$ in the index pair $(r,s)$ and then to the
	$\{\hat{L}^{\omega}_{a,b}\}$. Consequently the observables
	$\mathcal{O}^{\omega}
	= \sum_{a,b \in \F_2} (-1)^{a+b} \hat{L}^{\omega}_{a,b}$
	satisfy
	$\mathcal{O}^{x,z,r,s} \mathcal{O}^{x,z,r',s'}
	\approx_\delta \mathcal{O}^{x,z,r+r',s+s'}$
	on average; identifying $\F_q \cong \F_2^t$ so that the pair $(r,s)$
	ranges over $\F_2^{2t}$, Theorem~\ref{thm:linearity} applied for each
	fixed $(x,z)$ yields exactly linear observables
	$\{\mathcal{L}^{x,z,r,s}\}$ that are $\delta$-close to the
	$\{\mathcal{O}^{x,z,r,s}\}$ on average over $x, z, r, s$; no further
	extension of the state is needed because $\ket{\hat\psi}$ is padded with
	sufficiently many ancilla qubits in the zero state. Define
	$\hat{Q}^{x,z}_{a,b}
	= \E_{r,s} (-1)^{\tr(ra+sb)} \mathcal{L}^{x,z,r,s}$
	for $a, b \in \F_q$: exact linearity makes each $\hat{Q}^{x,z}_{a,b}$ a
	projection, and the family sums to the identity because
	$\mathcal{L}^{x,z,0,0} = \Id$, so $\{\hat{Q}^{x,z}_{a,b}\}$ is a
	projective measurement recovering $\F_q$-valued outcomes from the binary
	data by Fourier averaging over $(r,s)$. Its self-consistency
	\eqref{eq:qld-q-self-cons} follows from Lemma~\ref{lem:avg-closeness}
	together with the self-consistency of the $\{\mathcal{O}^{\omega}\}$;
	the consistency \eqref{eq:qld-q-cons-m-hat-zx} follows from a Fourier
	computation that replaces $\mathcal{L}$ by $\mathcal{O}$, then by
	$\hat{R}$, then by the ordered product of refined point measurements, and
	finally resolves the averages over $r$ and $s$ using
	Lemma~\ref{lem:cancellation}; and \eqref{eq:qld-q-cons-m-hat-xz} follows
	by inserting one application of the approximate commutation of
	Lemma~\ref{lem:qld-comm-cons} into the same chain. Repeating the argument
	with the tensor factors on registers $\mathsf{B}\mathsf{B}'$ and
	$\mathsf{A}\mathsf{B}''$ gives all symmetric equivalents.
\end{proof}

\section{Combined line measurements}

\begin{lemma}[Combined line measurements]
	\label{lem:qld-xz-lines}
	\uses{lem:qld-4-10, def:line, def:line-point-dist, def:consistency,
		def:bracket, def:expanded-state, def:symmetric-equivalents}
	There exists a function $\delta_P(\eps,m,d,q) = \poly(\eps, md/q)$ such
	that the following holds. For
	$\hilb \in \{\hilb_{\mathsf{A}}, \hilb_{\mathsf{B}}\}$ and every pair of
	lines $\ell_X$ and $\ell_Z$ in $\F_q^m$ there exists a POVM
	$\{T^{\ell_X,\ell_Z}_{f_X,f_Z}\}$ acting on $\hilb \ot (\C^q)^{\ot n}$,
	with answers consisting of pairs of polynomials
	$(f_X, f_Z) \in \deg_{md}(\ell_X) \times \deg_{md}(\ell_Z)$, such that on
	average over pairs $(\ell_X, x)$ and $(\ell_Z, z)$ independently sampled
	from the line-point distribution
	(Definition~\ref{def:line-point-dist}),
	\[
		\big(T^{\ell_X,\ell_Z}
		_{[\mathrm{eval}_x(\cdot)=a,\ \mathrm{eval}_z(\cdot)=b]}
		\big)_{\mathsf{A}\mathsf{A}'}
		\simeq_{\delta_P}
		\big(\hat{Q}^{x,z}_{a,b}\big)_{\mathsf{B}\mathsf{A}''}\,,
	\]
	as well as all symmetric equivalents of this approximation. Moreover, if
	$\ell_W$ is axis-parallel, then $f_W \in \deg_{d}(\ell_W)$.
\end{lemma}
\begin{proof}
	\uses{def:expanded-line-measurement, lem:qld-4-10, fact:add-a-proj,
		lem:cool-closeness-fact, lem:qld-comm-line-cons, fact:agreement,
		fact:data-processing, lem:pasting, lem:schwartz-zippel-total-degree}
	Define, for all lines $\ell_X, \ell_Z$ and
	$(f_X, f_Z) \in \deg_{md}(\ell_X) \times \deg_{md}(\ell_Z)$, the sandwich
	operator
	$T^{\ell_X,\ell_Z}_{f_X,f_Z}
	= \hat{M}^{(\mathrm{Line},X),\ell_X}_{f_X}
	\hat{M}^{(\mathrm{Line},Z),\ell_Z}_{f_Z}
	\hat{M}^{(\mathrm{Line},X),\ell_X}_{f_X}$,
	which forms a valid POVM; since the lines measurements come from a valid
	strategy for the low-degree test, on an axis-parallel line $\ell_W$ the
	outcome $f_W$ lies in $\deg_d(\ell_W)$, and this property is inherited by
	$T$. As intermediate consistency relations, write
	$\hat{Q}^{x,z}_{a,\bullet} = \sum_b \hat{Q}^{x,z}_{a,b}$ and
	$\hat{Q}^{x,z}_{\bullet,b} = \sum_a \hat{Q}^{x,z}_{a,b}$ for the two
	marginals of the combined point measurement of
	Lemma~\ref{lem:qld-4-10}. Applying
	Lemma~\ref{fact:add-a-proj} to \eqref{eq:qld-q-cons-m-hat-zx} and then
	Lemma~\ref{lem:cool-closeness-fact} gives, on average over uniform
	$x, z$,
	$(\hat{Q}^{x,z}_{a,\bullet})_{\mathsf{A}\mathsf{A}'}
	\approx_\delta
	\sum_b (\hat{Q}^{x,z}_{a,b})_{\mathsf{A}\mathsf{A}'}
	(\hat{M}^{(\mathrm{Point},Z),z}_{b})_{\mathsf{B}\mathsf{A}''}$,
	and a direct norm computation from this display and
	\eqref{eq:qld-q-cons-m-hat-zx} yields the marginal consistencies
	$(\hat{Q}^{x,z}_{a,\bullet})_{\mathsf{A}\mathsf{A}'} \approx_{\delta_Q}
	(\hat{M}^{(\mathrm{Point},X),x}_{a})_{\mathsf{B}\mathsf{A}''}$ and
	$(\hat{Q}^{x,z}_{\bullet,b})_{\mathsf{A}\mathsf{A}'} \approx_{\delta_Q}
	(\hat{M}^{(\mathrm{Point},Z),z}_{b})_{\mathsf{B}\mathsf{A}''}$.
	From items 1 and 3 of Lemma~\ref{lem:qld-comm-line-cons}, converting
	between $\approx$ and $\simeq$ with Lemma~\ref{fact:agreement} and
	summing outcomes with Lemma~\ref{fact:data-processing}, the lines
	measurements are consistent with the points measurements, and combining
	with the marginal consistencies gives, on average over the line-point
	distribution,
	$(\hat{Q}^{x,z}_{a,\bullet})_{\mathsf{A}\mathsf{A}'} \simeq_{\delta_Q}
	(\hat{M}^{(\mathrm{Line},X),\ell_X}
	_{[\mathrm{eval}_x(\cdot)=a]})_{\mathsf{B}\mathsf{A}''}$
	together with the analogous relation for $Z$. The conclusion now follows
	from the pasting lemma (Lemma~\ref{lem:pasting}) applied with
	$\hat{Q}^{x,z}_{a,b}$ in the role of the combined measurement,
	$\hat{M}^{(\mathrm{Line},X),\ell_X}$ and
	$\hat{M}^{(\mathrm{Line},Z),\ell_Z}$ in the roles of the two factor
	measurements, $T^{\ell_X,\ell_Z}$ in the role of the pasted measurement,
	$\delta := \delta_Q$, and $\eta := md/q$, the last justified by the
	Schwartz--Zippel bound (Lemma~\ref{lem:schwartz-zippel-total-degree})
	applied to the polynomials $f_X, f_Z$ of total degree at most $md$; the
	two hypotheses of the pasting lemma are supplied by the preceding display
	and by \eqref{eq:qld-q-self-cons}, and the distribution over
	$(\ell_X, \ell_Z)$ is that of two independent uniformly random lines.
	This gives the stated relation with $\delta_P$ equal to the pasting
	error, a $\poly(\eps, md/q)$ function, and the identical argument with
	the registers modified gives all symmetric equivalents.
\end{proof}

\section{Applying the classical low-degree test}
\label{sec:apply-ldt}

The combined point and line measurements of the previous two sections are next
merged into a single measurement returning a pair of global polynomials that
encode the $X$-basis and the $Z$-basis information simultaneously. This is
done by encoding the pair as one polynomial in $2m+2$ variables and invoking
the quantum soundness of the classical low individual degree test
(Theorem~\ref{lem:ld-soundness}). The instantiation at dimension $2m+2$
requires that $2m+2$ divide $q$; the statements below carry this hypothesis
explicitly, and Remark~\ref{rem:qld-4-7-divisibility} discusses its
relationship to the admissibility of $(q,m,d)$.

\begin{definition}[The combining map]
	\label{def:combine-map}
	\uses{def:polyfunc, def:line}
	For $f, g \in \mathrm{ideg}_{d,m}(\F_q)$, the polynomial
	$\mathrm{combine}_{f,g} \in \mathrm{ideg}_{d,2m+2}(\F_q)$ is defined by
	\[
		\mathrm{combine}_{f,g}(x, z, \alpha, \beta)
		= \alpha f(x) + \beta g(z)\,,
		\qquad x \in \F_q^m,\ z \in \F_q^m,\ \alpha \in \F_q,\
		\beta \in \F_q\,.
	\]
	The combining map is also defined for polynomials restricted to lines.
	Let $\ell$ be a line in $\F_q^{2m+2}$ with intercept
	$u = (u_X, u_Z, u_\alpha, u_\beta)$ and slope
	$v = (v_X, v_Z, v_\alpha, v_\beta)$, and let $\ell_X$ and $\ell_Z$ be
	lines in $\F_q^m$ such that for every point
	$(x, z, \alpha, \beta) \in \ell$ it holds that $x \in \ell_X$ and
	$z \in \ell_Z$. Then for any two polynomials
	$f \in \deg_{md}(\ell_X)$ and $g \in \deg_{md}(\ell_Z)$, the polynomial
	$\mathrm{combine}_{f,g} \in \deg_{md+1}(\ell)$ is defined by
	\begin{equation}
		\label{eq:combine-lines}
		\mathrm{combine}_{f,g}(u + t \cdot v)
		= (u_\alpha + t \cdot v_\alpha)\, f(u_X + t \cdot v_X)
		+ (u_\beta + t \cdot v_\beta)\, g(u_Z + t \cdot v_Z)\,,
		\qquad t \in \F_q\,.
	\end{equation}
	The degree bound holds because $f$ and $g$, restricted to lines, are
	polynomials of degree at most $md$ in the line parameter, and each is
	multiplied by an affine factor of degree at most one.
\end{definition}

\begin{lemma}[Combined point measurement on the extended coordinates]
	\label{lem:qld-4-12}
	\uses{lem:qld-4-10, def:expanded-state, def:expanded-point-measurement,
		def:povm-distance, def:symmetric-equivalents}
	For all $x, z \in \F_q^m$, $\alpha, \beta \in \F_q$, and
	$\hilb \in \{\hilb_{\mathsf{A}}, \hilb_{\mathsf{B}}\}$ there exists a
	projective measurement
	$\{\hat{Q}^{x,z,\alpha,\beta}_{c}\}_{c \in \F_q}$ acting on
	$\hilb \ot (\C^q)^{\ot n}$ such that the following holds, where
	$\delta_Q$ is the function of Lemma~\ref{lem:qld-4-10}:
	\begin{enumerate}
		\item (\textbf{Self-consistency}) On average over uniformly random
		$(x, z, \alpha, \beta) \in \F_q^{2m+2}$,
		\begin{equation}
			\label{eq:qld-4-12-self-cons}
			(\hat{Q}^{x,z,\alpha,\beta}_{c})_{\mathsf{A}\mathsf{A}'}
			\approx_{\delta_Q}
			(\hat{Q}^{x,z,\alpha,\beta}_{c})_{\mathsf{B}\mathsf{A}''}\,.
		\end{equation}
		\item (\textbf{Consistency with $\hat{M}$}) On average over uniformly
		random $(x, z, \alpha, \beta) \in \F_q^{2m+2}$,
		\begin{gather}
			\label{eq:qld-4-12-cons-m-hat}
			(\hat{Q}^{x,z,\alpha,\beta}_{c})_{\mathsf{A}\mathsf{A}'}
			\approx_{\delta_Q}
			\Big( \sum_{\substack{a,b:\\ \alpha a + \beta b = c}}
			\hat{M}^{(\mathrm{Point},X),x}_{a} \,
			\hat{M}^{(\mathrm{Point},Z),z}_{b}
			\Big)_{\mathsf{B}\mathsf{A}''}\,, \\
			\label{eq:qld-4-12-cons-m-hat2}
			(\hat{Q}^{x,z,\alpha,\beta}_{c})_{\mathsf{A}\mathsf{A}'}
			\approx_{\delta_Q}
			\Big( \sum_{\substack{a,b:\\ \alpha a + \beta b = c}}
			\hat{M}^{(\mathrm{Point},Z),z}_{b} \,
			\hat{M}^{(\mathrm{Point},X),x}_{a}
			\Big)_{\mathsf{B}\mathsf{A}''}\,.
		\end{gather}
	\end{enumerate}
	Furthermore, all symmetric equivalents of these approximations also hold.
\end{lemma}
\begin{proof}
	\uses{lem:qld-4-10, fact:data-processing}
	Define
	$\hat{Q}^{x,z,\alpha,\beta}_{c}
	= \sum_{a,b \,:\, \alpha a + \beta b = c} \hat{Q}^{x,z}_{a,b}$,
	a coarse-graining of the projective measurement of
	Lemma~\ref{lem:qld-4-10}, hence itself projective. A uniformly random
	point of $\F_q^{2m+2}$ has uniformly random coordinates $(x,z)$, so the
	self-consistency and consistency properties follow from those of
	$\{\hat{Q}^{x,z}_{a,b}\}$ established in Lemma~\ref{lem:qld-4-10}
	together with Lemma~\ref{fact:data-processing}, which shows that
	coarse-graining the outcomes preserves the $\delta_Q$ bounds.
\end{proof}

\begin{lemma}[Combined line measurement on the extended coordinates]
	\label{lem:qld-4-13}
	\uses{lem:qld-4-12, def:line, def:line-point-dist, def:consistency,
		def:bracket, def:expanded-state}
	% The source states delta_combine(eps,m,d,q) = poly(m^2 eps, md/q)
	% (references/qpbt-paper/14_analysis_of_the_pauli_basis_test.tex,
	% lem:qld-4-13); neither of its two proof routes delivers that form,
	% and the route retained here yields m * poly(eps, md/q). See
	% rem:qld-4-13-source-defects. The divisibility hypothesis is required
	% for the line-point distribution over F_q^{2m+2}; see
	% rem:qld-4-7-divisibility.
	There exists a function
	$\delta_{\mathrm{combine}}(\eps,m,d,q) = m \cdot \poly(\eps, md/q)$
	such that the following holds. Assume that $2m+2$ divides $q$. For
	every line $\ell$ in $\F_q^{2m+2}$ and
	$\hilb \in \{\hilb_{\mathsf{A}}, \hilb_{\mathsf{B}}\}$ there exists a
	POVM $\{\hat{Q}^{\ell}_{f}\}$ acting on $\hilb \ot (\C^q)^{\ot n}$ with
	outcomes $f \in \deg_{md+1}(\ell)$ such that on average over a uniformly
	random pair $(\ell, (x,z,\alpha,\beta))$ drawn from the line-point
	distribution over $\F_q^{2m+2}$,
	\begin{gather}
		\label{eq:qld-4-13}
		\big(\hat{Q}^{\ell}
		_{[\mathrm{eval}_{(x,z,\alpha,\beta)}(\cdot)=a]}
		\big)_{\mathsf{A}\mathsf{A}'}
		\simeq_{\delta_{\mathrm{combine}}}
		\big(\hat{Q}^{x,z,\alpha,\beta}_{a}\big)_{\mathsf{B}\mathsf{A}''}\,,
		\\
		\big(\hat{Q}^{\ell}
		_{[\mathrm{eval}_{(x,z,\alpha,\beta)}(\cdot)=a]}
		\big)_{\mathsf{B}\mathsf{B}'}
		\simeq_{\delta_{\mathrm{combine}}}
		\big(\hat{Q}^{x,z,\alpha,\beta}_{a}\big)_{\mathsf{A}\mathsf{B}''}\,,
		\notag
	\end{gather}
	where the answer summation is over $a \in \F_q$. Moreover, if $\ell$ is
	axis-parallel, then the outcome $f \in \deg_{d}(\ell)$.
\end{lemma}

Before proving the lemma we record definitions and estimates concerning
distributions over lines and points. Recall from
Definition~\ref{def:line-point-dist} that the line-point distribution over
$\F_q^m$ is an even mixture of an axis-parallel component $D_{\mathrm{ALine}}$
and a diagonal component $D_{\mathrm{DLine}}$.

\begin{definition}[Restricted line distributions]
	\label{def:ith-restricted-line}
	\uses{def:line-point-dist, def:ld-question-distribution}
	For any $i \in \{1, \dots, m\}$, the \emph{$i$-th restricted
	axis-parallel line distribution} $D_{\mathrm{ALine},i}$ is the
	restriction of $D_{\mathrm{ALine}}$ to pairs $(\ell, u)$ where the line
	$\ell$ is along the direction $e_i$, i.e.\ is of the form
	$\ell = (u_0, e_i)$. Similarly, the \emph{$i$-th restricted diagonal line
	distribution} $D_{\mathrm{DLine},i}$ is the restriction of
	$D_{\mathrm{DLine}}$ to pairs $(\ell, u)$ where the line
	$\ell = (u_0, s, v)$ satisfies $\chi(s) = i$ and coordinates
	$1, \dots, i-1$ of $v$ are $0$. (The function $\chi$ is the one used in
	the question distribution of the low individual degree game,
	Definition~\ref{def:ld-question-distribution}.)
\end{definition}

\begin{lemma}[Restricted line distributions and error inflation]
	\label{lem:restricted-line-mixture-bounds}
	% Source: unlabelled prose in the paper --- the discussion following
	% def:ith-restricted-line at lines 1049-1061 of
	% references/qpbt-paper/14_analysis_of_the_pauli_basis_test.tex.
	\uses{def:ith-restricted-line, def:line-point-dist, lem:qld-xz-lines,
		lem:qld-4-10, def:expanded-state}
	The distribution $D_{\mathrm{ALine}}$ is a uniform mixture of the
	distributions $D_{\mathrm{ALine},i}$ for $i \in \{1, \dots, m\}$, and
	likewise $D_{\mathrm{DLine}}$ is a uniform mixture of the
	$D_{\mathrm{DLine},i}$. As a consequence:
	\begin{enumerate}
		\item any approximation that holds on average over the line-point
		distribution $D_{\mathrm{Line}}$ with error $\delta$ holds over each
		restricted distribution $D_{\mathrm{ALine},i}$ and
		$D_{\mathrm{DLine},i}$ with error $2m\delta$;
		\item any approximation that holds on average over
		$D_{\mathrm{Line}} \times D_{\mathrm{Line}}$ holds over any product
		of restricted distributions with error $4m^2\delta$;
		\item in particular, by Lemma~\ref{lem:qld-xz-lines}, for any types
		$\mathrm{v}_1, \mathrm{v}_2 \in \{\mathrm{ALine},
		\mathrm{DLine}\}$ and any $i, j \in \{1, \dots, m\}$,
		\begin{equation}
			\label{eq:qld-xz-lines-restricted}
			\E_{(\ell_X,x) \sim D_{\mathrm{v}_1,i}}\
			\E_{(\ell_Z,z) \sim D_{\mathrm{v}_2,j}}
			\sum_{f_X,f_Z}
			\bra{\hat\psi}
			\big(T^{\ell_X,\ell_Z}_{f_X,f_Z}\big)_{\mathsf{A}\mathsf{A}'}
			\ot
			\big(\Id - \hat{Q}^{x,z}_{f_X(x),f_Z(z)}
			\big)_{\mathsf{B}\mathsf{A}''}
			\ket{\hat\psi}
			= O(m^2 \delta_P)\,.
		\end{equation}
	\end{enumerate}
\end{lemma}
\begin{proof}
	\uses{lem:qld-xz-lines, def:line-point-dist, def:consistency}
	The line-point distribution places weight $\tfrac{1}{2}$ on each of the
	axis-parallel and diagonal components and, within each component, uniform
	weight on the index $i$, so each restricted distribution carries mixture
	weight $\tfrac{1}{2m}$ in $D_{\mathrm{Line}}$. A nonnegative average that
	is at most $\delta$ therefore bounds each mixture component by
	$2m\delta$, and products of two such mixtures give the factor
	$(2m)^2 = 4m^2$. The instance \eqref{eq:qld-xz-lines-restricted} is item
	2 applied to the conclusion of Lemma~\ref{lem:qld-xz-lines}, whose
	$\simeq_{\delta_P}$ relation over two independent copies of the
	line-point distribution is precisely the statement that the left-hand
	side of \eqref{eq:qld-xz-lines-restricted}, averaged over
	$D_{\mathrm{Line}} \times D_{\mathrm{Line}}$, is at most $\delta_P$.
\end{proof}

\begin{lemma}[The sub-line distribution]
	\label{lem:qld-sublines}
	\uses{def:line, def:line-point-dist, def:ith-restricted-line}
	% Hypothesis added: the line-point distribution over F_q^{2m+2} is built
	% from the index function chi of def:ld-question-distribution at
	% dimension 2m+2, which requires 2m+2 to divide q; the source
	% (references/qpbt-paper/14_analysis_of_the_pauli_basis_test.tex,
	% lem:qld-sublines) instantiates it without this hypothesis. See
	% rem:qld-4-7-divisibility.
	Assume that $2m+2$ divides $q$. Then there exists a distribution $D$
	over tuples $(\ell, \ell_X, \ell_Z)$, where $\ell$ is a line in
	$\F_q^{2m+2}$ and $\ell_X, \ell_Z$ are lines in $\F_q^m$, with the
	following properties:
	\begin{enumerate}
		\item The marginal distribution over $\ell$ is the same as the
		marginal of the line-point distribution over $\F_q^{2m+2}$.
		% The source phrases the next property as: the marginal over
		% (ell_X, ell_Z) is "a mixture of the product distributions
		% D_{ALine,i} x D_{ALine,j} and D_{DLine,i} x D_{DLine,j}". Those
		% products are distributions over pairs of line-point pairs, not
		% over pairs of lines, so the phrase does not type-check as
		% written; the formulation below states the mixture at the level
		% of the joint law of (ell_X, ell_Z, u), which is what the proofs
		% of lem:claim-17-2 and lem:claim-17-3 use.
		\item For any point of the form $u = (u_X, u_Z, \alpha, \beta)$
		lying on $\ell$, it holds that $u_X \in \ell_X$ and
		$u_Z \in \ell_Z$. Moreover, the joint distribution of
		$(\ell_X, \ell_Z, u)$, for $(\ell, \ell_X, \ell_Z) \sim D$ and
		$u$ chosen uniformly at random on $\ell$, is a mixture of
		components indexed by a type
		$\mathrm{v} \in \{\mathrm{ALine}, \mathrm{DLine}\}$ and indices
		$i, j \in \{1, \dots, m\}$ with the following two properties
		within the component $(\mathrm{v}, i, j)$: the pair
		$(\ell_X, \ell_Z)$ is distributed according to the product of the
		marginal distributions on lines of $D_{\mathrm{v},i}$ and of
		$D_{\mathrm{v},j}$, and, conditioned on $(\ell_X, \ell_Z)$, the
		point $u_X$ is uniformly distributed on $\ell_X$ and likewise the
		point $u_Z$ is uniformly distributed on $\ell_Z$. Equivalently,
		since under each restricted distribution the point is uniformly
		distributed on the line conditionally on the line: within the
		component $(\mathrm{v}, i, j)$, the pair
		$((\ell_X, u_X), (\ell_Z, z'))$, where $z'$ is a uniformly random
		point of $\ell_Z$ sampled independently of $u$, is distributed
		according to $D_{\mathrm{v},i} \times D_{\mathrm{v},j}$, and
		symmetrically the pair $((\ell_X, x'), (\ell_Z, u_Z))$, with $x'$
		a uniformly random point of $\ell_X$ sampled independently of
		$u$, is distributed according to
		$D_{\mathrm{v},i} \times D_{\mathrm{v},j}$. (No assertion is made
		about the joint conditional law of the pair $(u_X, u_Z)$.)
		\item If $\ell$ is axis-parallel, then $\ell_X$ and $\ell_Z$ are
		axis-parallel.
	\end{enumerate}
	(Property~3 is not part of the numbered statement in the source; see
	Remark~\ref{rem:qld-sublines-property-three}.)
\end{lemma}
\begin{proof}
	\uses{def:line-representative, def:ld-question-distribution,
		def:line-point-dist, def:ith-restricted-line, def:line,
		prop:line-equiv}
	The distribution $D$ is described by a sampling procedure. First sample
	$(\ell, u)$ from the line-point distribution over $\F_q^{2m+2}$ and write
	$u = (u_X, u_Z, \alpha, \beta)$, so that Property~1 holds by
	construction; all auxiliary randomness introduced below is sampled
	independently of $(\ell, u)$. We write $\chi$ for the index function of
	Definition~\ref{def:ld-question-distribution} at dimension $m$, whose
	classes $\chi^{-1}(i)$ for $i \in \{1, \dots, m\}$ partition $\F_q$
	into sets of equal size $q/m$, and $\chi_{2m+2}$ for the index function
	at dimension $2m+2$, defined because $2m+2$ divides $q$. By a
	\emph{fresh diagonal direction pair} we mean a pair $(s', w')$ with
	$s'$ uniformly random in $\F_q$ and $w' \in \F_q^m$ having coordinates
	$1, \dots, \chi(s')-1$ equal to $0$ and the remaining coordinates
	uniformly random; this is exactly the law of the coordinate--direction
	data of the diagonal component of the line-point distribution over
	$\F_q^m$ (Definition~\ref{def:line-point-dist}).

	If $\ell$ is axis-parallel, $\ell = (v, e_j)$ with
	$v = (v_X, v_Z, \alpha_0, \beta_0)$ and $j \in \{1, \dots, 2m+2\}$, set
	$\ell_X = (v'_X, e_{i_X})$ and $\ell_Z = (v'_Z, e_{i_Z})$ where: for
	$1 \le j \le m$, $i_X = j$ and $i_Z$ is uniformly random in
	$\{1,\dots,m\}$; for $m+1 \le j \le 2m$, $i_X$ is uniformly random and
	$i_Z = j - m$; for $2m+1 \le j \le 2m+2$, both $i_X$ and $i_Z$ are
	uniformly random; and $v'_X, v'_Z$ are the canonical line
	representatives of $v_X, v_Z$ for the chosen directions
	(Definition~\ref{def:line-representative}), so that $v_X \in \ell_X$
	and $v_Z \in \ell_Z$.

	If $\ell$ is diagonal, $\ell = (v, s, w)$, then by the construction of
	the diagonal component of the line-point distribution
	(Definition~\ref{def:ld-question-distribution} instantiated at
	dimension $2m+2$) the direction $w = (w_X, w_Z, \gamma, \delta)$ has
	coordinates $1, \dots, j-1$ equal to $0$, where $j = \chi_{2m+2}(s)$,
	and its remaining coordinates are uniformly random. These remaining
	coordinates may also vanish, so $w$ need not have a nonzero coordinate
	at index $j$ and may be the zero vector, in which case $\ell$ is a
	singleton (Definition~\ref{def:line}).
	% The source proof describes w as "a direction vector whose first
	% nonzero coordinate is j = chi(s)"
	% (references/qpbt-paper/14_analysis_of_the_pauli_basis_test.tex, proof
	% of lem:qld-sublines); the distribution delivers only the weaker
	% property recorded here, which is the one the argument uses. The
	% source also leaves the laws of s_X and s_Z unspecified ("chosen such
	% that chi(s_X) = j"); the uniform choices below are needed for the
	% exact marginals claimed in Property 2.
	Sample $\ell_X = (v'_X, s_X, w'_X)$ and $\ell_Z = (v'_Z, s_Z, w'_Z)$
	where: for $1 \le j \le m$, $w'_X = w_X$ and $w'_Z = w_Z$, with $s_X$
	uniformly random in the class $\chi^{-1}(j)$ and $s_Z$ uniformly
	random in the class $\chi^{-1}(1)$; for $m+1 \le j \le 2m$, the pair
	$(s_X, w'_X)$ is a fresh diagonal direction pair and $w'_Z = w_Z$,
	with $s_Z$ uniformly random in the class $\chi^{-1}(j-m)$; for
	$2m+1 \le j \le 2m+2$, both $(s_X, w'_X)$ and $(s_Z, w'_Z)$ are fresh
	diagonal direction pairs, sampled independently; and $v'_X, v'_Z$ are
	the canonical representatives of $v_X, v_Z$ for the directions
	$w'_X, w'_Z$.

	Property~3 is immediate from the procedure: when $\ell$ is
	axis-parallel, both $\ell_X$ and $\ell_Z$ are constructed to be
	axis-parallel. The containments in Property~2 hold for every point
	$u^\dagger = (u^\dagger_X, u^\dagger_Z, \alpha^\dagger,
	\beta^\dagger)$ of $\ell$: the block $u^\dagger_X$ differs from $v_X$
	by a multiple of the direction of $\ell_X$ when that direction is
	inherited from $\ell$ (namely $e_j$ for axis-parallel $\ell$ with
	$1 \le j \le m$, and $w_X$ for diagonal $\ell$ with $1 \le j \le m$),
	and equals $v_X$ in all other cases, since then the corresponding
	block of the direction of $\ell$ vanishes; in either case
	$u^\dagger_X \in \ell_X$ because $v_X \in \ell_X$, and symmetrically
	$u^\dagger_Z \in \ell_Z$.

	For the distributional part of Property~2, observe first that
	$(\ell_X, \ell_Z)$ is a function of $\ell$ and of the auxiliary
	randomness alone: the directions used are either read off from the
	direction data of $\ell$ or freshly sampled, and by
	Proposition~\ref{prop:line-equiv} the canonical representatives
	$v'_X, v'_Z$ do not depend on which base point of $\ell$ is used to
	compute them. Under both components of the line-point distribution the
	point $u$ is uniformly distributed on $\F_q^{2m+2}$ and, conditioned
	on the line $\ell$, uniformly distributed on the points of $\ell$
	(each of the $q$ values of the line parameter, or the single point
	when the direction vanishes, being equally likely). Hence, conditioned
	on $(\ell, \ell_X, \ell_Z)$, the sampled point $u$ is uniformly
	distributed on $\ell$, and it suffices to establish the claimed
	mixture description for the tuple $(\ell_X, \ell_Z, u)$ produced by
	the procedure itself.

	Condition on the component of the line-point distribution
	(axis-parallel or diagonal) and on $j$; the mixture of Property~2 is
	over these choices together with the fresh index choices, and the
	blocks $u_X, u_Z$ of $u$ are independent and uniform on $\F_q^m$,
	independent of all direction data. In the axis-parallel case with
	$1 \le j \le m$, the pair $(\ell_X, u_X)$ consists of the uniformly
	random point $u_X$ together with the canonical axis-parallel line
	through it in direction $e_j$, so
	$(\ell_X, u_X) \sim D_{\mathrm{ALine},j}$, while, conditioned
	additionally on $i_Z$, the pair $(\ell_Z, u_Z)$ consists of the
	uniformly random point $u_Z$ together with the canonical line through
	it in direction $e_{i_Z}$, so
	$(\ell_Z, u_Z) \sim D_{\mathrm{ALine},i_Z}$; the two pairs are
	functions of the disjoint independent data $u_X$ and $(u_Z, i_Z)$,
	hence independent. In the diagonal case with $1 \le j \le m$, the
	triple $(u_X, w_X, s_X)$ has $u_X$ uniform, $w_X$ with coordinates
	$1, \dots, j-1$ equal to $0$ and the remaining coordinates uniform,
	and $s_X$ uniform in $\chi^{-1}(j)$, which is exactly the law
	generating $D_{\mathrm{DLine},j}$; hence
	$(\ell_X, u_X) \sim D_{\mathrm{DLine},j}$, and similarly
	$(\ell_Z, u_Z) \sim D_{\mathrm{DLine},1}$ since $w_Z$ is uniform and
	$s_Z$ is uniform in $\chi^{-1}(1)$; the two pairs are functions of the
	disjoint independent data $(u_X, w_X, s_X)$ and $(u_Z, w_Z, s_Z)$,
	hence independent. The remaining cases are analogous: each of
	$(\ell_X, u_X)$ and $(\ell_Z, u_Z)$ is either built from the inherited
	direction block as above, or from a fresh axis index or fresh diagonal
	direction pair, in which case it is distributed as the uniform mixture
	over $i$ of $D_{\mathrm{v},i}$ for the corresponding type
	$\mathrm{v}$; and the two pairs are always functions of disjoint
	independent collections of the variables $u_X$, $u_Z$, the direction
	blocks, and the fresh auxiliary choices, hence independent conditioned
	on the component. Within each component this yields exactly the
	description of Property~2: the pair $(\ell_X, \ell_Z)$ is
	product-distributed according to the line marginals of
	$D_{\mathrm{v},i}$ and $D_{\mathrm{v},j}$, and, conditioned on
	$(\ell_X, \ell_Z)$, each of $u_X$ and $u_Z$ is uniformly distributed
	on its line, because under each restricted distribution the point is
	uniform on the line conditionally on the line. The equivalent product
	formulations follow from the same conditional uniformity.
\end{proof}

\begin{remark}[Property~3 of the sub-line distribution]
	\label{rem:qld-sublines-property-three}
	The source proof of Lemma~\ref{lem:qld-4-13} in \cite{Ji2020MIPStar}
	appeals to a ``Property~3'' of the sub-line lemma, while the statement
	given there lists only two numbered properties, so the reference is
	dangling. The intended assertion is the axis-parallel closure property
	recorded as Property~3 of Lemma~\ref{lem:qld-sublines} above, which is
	immediate from the sampling procedure in the proof.
\end{remark}

The proof of Lemma~\ref{lem:qld-4-13} rests on three scalar estimates, stated
next. In all three, $2m+2$ is assumed to divide $q$, $D$ is the distribution
of Lemma~\ref{lem:qld-sublines},
$T^{\ell_X,\ell_Z}$ is the measurement of Lemma~\ref{lem:qld-xz-lines} placed
on registers $\mathsf{A}\mathsf{A}'$, the second tensor factor is placed on
registers $\mathsf{B}\mathsf{A}''$, and the sums range over
$(f_X, f_Z) \in \deg_{md}(\ell_X) \times \deg_{md}(\ell_Z)$.

\begin{lemma}[Replacing the combined point measurement by the ordered product]
	\label{lem:claim-17-1}
	% claim:17-1 in the paper; the original label is recorded here as a
	% comment for cross-reference fidelity.
	\uses{lem:qld-sublines, lem:qld-xz-lines, lem:qld-4-10,
		def:expanded-point-measurement, def:expanded-state}
	\begin{align*}
		\E_{(\ell,\ell_X,\ell_Z) \sim D}\
		\E_{(x,z,\alpha,\beta) \in \ell}
		&\sum_{f_X,f_Z} \bra{\hat\psi}
		\big(T^{\ell_X,\ell_Z}_{f_X,f_Z}\big)_{\mathsf{A}\mathsf{A}'} \ot
		\big(\hat{Q}^{x,z}_{f_X(x),f_Z(z)}\big)_{\mathsf{B}\mathsf{A}''}
		\ket{\hat\psi} \\
		\approx_{\delta_Q^{1/2}}\
		\E_{(\ell,\ell_X,\ell_Z) \sim D}\
		\E_{(x,z,\alpha,\beta) \in \ell}
		&\sum_{f_X,f_Z} \bra{\hat\psi}
		\big(T^{\ell_X,\ell_Z}_{f_X,f_Z}\big)_{\mathsf{A}\mathsf{A}'} \ot
		\big(\hat{M}^{(\mathrm{Point},Z),z}_{f_Z(z)} \,
		\hat{M}^{(\mathrm{Point},X),x}_{f_X(x)}
		\big)_{\mathsf{B}\mathsf{A}''}
		\ket{\hat\psi}\,.
	\end{align*}
\end{lemma}
\begin{proof}
	\uses{lem:qld-4-10, lem:qld-sublines, lem:qld-xz-lines,
		def:line-point-dist}
	Write
	$W^{x,z}_{a,b} = \hat{Q}^{x,z}_{a,b}
	- \hat{M}^{(\mathrm{Point},Z),z}_{b}
	\hat{M}^{(\mathrm{Point},X),x}_{a}$
	and apply the Cauchy--Schwarz inequality to the difference of the two
	sides. The first factor is at most $1$ because
	$\sum_{f_X,f_Z \,:\, f_X(x)=a,\ f_Z(z)=b}
	T^{\ell_X,\ell_Z}_{f_X,f_Z} \le \Id$
	for every $x, z, a, b$. For the second factor, note that when $\ell$ is
	drawn from the line-point distribution over $\F_q^{2m+2}$, a uniformly
	random point $(x,z,\alpha,\beta)$ on $\ell$ is a uniformly random point
	of $\F_q^{2m+2}$, so the coordinates $x, z$ are independent and uniform
	in $\F_q^m$; hence
	$\E_{x,z} \sum_{a,b}
	\bra{\hat\psi} \Id \ot (W^{x,z}_{a,b})^2 \ket{\hat\psi}
	\le \delta_Q$
	by \eqref{eq:qld-q-cons-m-hat-zx}. The total error is therefore at most
	$\sqrt{\delta_Q}$.
\end{proof}

\begin{lemma}[Removing the $X$-point factor]
	\label{lem:claim-17-2}
	% claim:17-2 in the paper; the original label is recorded here as a
	% comment for cross-reference fidelity. The source right-hand side
	% carries a typographical superscript "Z" in place of the point "z";
	% the corrected form is stated here.
	\uses{lem:qld-sublines, lem:qld-xz-lines,
		def:expanded-point-measurement, lem:qld-comm-line-cons,
		def:expanded-state}
	\begin{align*}
		\E_{(\ell,\ell_X,\ell_Z) \sim D}\
		\E_{(x,z,\alpha,\beta) \in \ell}
		&\sum_{f_X,f_Z} \bra{\hat\psi}
		\big(T^{\ell_X,\ell_Z}_{f_X,f_Z}\big)_{\mathsf{A}\mathsf{A}'} \ot
		\big(\hat{M}^{(\mathrm{Point},Z),z}_{f_Z(z)} \,
		\hat{M}^{(\mathrm{Point},X),x}_{f_X(x)}
		\big)_{\mathsf{B}\mathsf{A}''}
		\ket{\hat\psi} \\
		\approx_{m\,\delta_{\mathrm{Line}}^{1/2}}\
		\E_{(\ell,\ell_X,\ell_Z) \sim D}\
		\E_{(x,z,\alpha,\beta) \in \ell}
		&\sum_{f_X,f_Z} \bra{\hat\psi}
		\big(T^{\ell_X,\ell_Z}_{f_X,f_Z}\big)_{\mathsf{A}\mathsf{A}'} \ot
		\big(\hat{M}^{(\mathrm{Point},Z),z}_{f_Z(z)}
		\big)_{\mathsf{B}\mathsf{A}''}
		\ket{\hat\psi}\,,
	\end{align*}
	where $\delta_{\mathrm{Line}}$ is the error function of
	Lemma~\ref{lem:qld-comm-line-cons}.
\end{lemma}
\begin{proof}
	\uses{lem:qld-comm-line-cons, lem:qld-sublines, lem:qld-xz-lines,
		lem:restricted-line-mixture-bounds, def:ith-restricted-line,
		def:expanded-line-measurement}
	By Cauchy--Schwarz, the difference of the two sides is bounded by the
	square root of the deficit
	\[
		1 - \E_{(\ell,\ell_X,\ell_Z) \sim D}\
		\E_{(x,z,\alpha,\beta) \in \ell} \sum_{a}
		\bra{\hat\psi}
		T^{\ell_X,\ell_Z,x}_{a} \ot
		\hat{M}^{(\mathrm{Point},X),x}_{a}
		\ket{\hat\psi}\,,
		\qquad
		T^{\ell_X,\ell_Z,x}_{a}
		= \sum_{f_X,f_Z \,:\, f_X(x)=a}
		T^{\ell_X,\ell_Z}_{f_X,f_Z}\,.
	\]
	By the sandwich definition of $T$ and the projectivity of the lines
	measurement $\{\hat{M}^{(\mathrm{Line},X),\ell_X}_{f_X}\}$, the
	subtracted expectation equals
	$\E \sum_a \bra{\hat\psi}
	\hat{M}^{(\mathrm{Line},X),\ell_X}_{[\mathrm{eval}_x(\cdot)=a]} \ot
	\hat{M}^{(\mathrm{Point},X),x}_{a} \ket{\hat\psi}$,
	which by Property~2 of Lemma~\ref{lem:qld-sublines} is a mixture of the
	corresponding averages over the restricted distributions
	$D_{\mathrm{v},i}$. By the error inflation of
	Lemma~\ref{lem:restricted-line-mixture-bounds} applied to the
	line--point consistency in item~3 of Lemma~\ref{lem:qld-comm-line-cons},
	each restricted term has deficit $O(m\,\delta_{\mathrm{Line}})$, so the
	total deficit is $O(m^2\,\delta_{\mathrm{Line}})$ (the final line of the
	source proof asserts this bound for the consistency term itself rather
	than for its deficit; see Remark~\ref{rem:qld-4-13-source-defects}).
	Taking the square root gives the claimed error
	$m\,\delta_{\mathrm{Line}}^{1/2}$.
\end{proof}

\begin{lemma}[The remaining $Z$-point term is close to one]
	\label{lem:claim-17-3}
	% claim:17-3 in the paper; the original label is recorded here as a
	% comment for cross-reference fidelity.
	\uses{lem:qld-sublines, lem:qld-xz-lines, lem:qld-4-10,
		def:expanded-point-measurement, def:expanded-state}
	\[
		\E_{(\ell,\ell_X,\ell_Z) \sim D}\
		\E_{(x,z,\alpha,\beta) \in \ell}
		\sum_{f_X,f_Z} \bra{\hat\psi}
		\big(T^{\ell_X,\ell_Z}_{f_X,f_Z}\big)_{\mathsf{A}\mathsf{A}'} \ot
		\big(\hat{M}^{(\mathrm{Point},Z),z}_{f_Z(z)}
		\big)_{\mathsf{B}\mathsf{A}''}
		\ket{\hat\psi}
		\ \approx_{\sqrt{m}\,(\delta_P^{1/4} + \delta_Q^{1/4} +
			\eps^{1/4})}\ 1\,.
	\]
\end{lemma}
\begin{proof}
	\uses{lem:restricted-line-mixture-bounds, lem:qld-sublines,
		lem:qld-xz-lines, lem:qld-4-10, lem:qld-comm-cons,
		def:ith-restricted-line}
	By Cauchy--Schwarz it suffices to bound the deficit
	$1 - \E \sum_b \bra{\hat\psi} T^{\ell_X,\ell_Z,z}_{b} \ot
	\hat{M}^{(\mathrm{Point},Z),z}_{b} \ket{\hat\psi}$,
	where
	$T^{\ell_X,\ell_Z,z}_{b}
	= \sum_{f_X} T^{\ell_X,\ell_Z}_{f_X,[\mathrm{eval}_z(\cdot)=b]}$.
	By Property~2 of Lemma~\ref{lem:qld-sublines}, the expectation can be
	rewritten over the mixture of product distributions
	$D_{\mathrm{v},i} \times D_{\mathrm{v},j}$ with $z$ uniformly random on
	$\ell_Z$, and, since
	$T^{\ell_X,\ell_Z,z}_{b}
	= \E_{x \in \ell_X} \sum_a
	T^{\ell_X,\ell_Z}
	_{[\mathrm{eval}_x(\cdot)=a,\ \mathrm{eval}_z(\cdot)=b]}$,
	one may additionally average over $x$ uniformly random on $\ell_X$. Up
	to error $m(\sqrt{\delta_P} + \sqrt{\delta_Q})$ --- obtained by
	applying the error inflation of
	Lemma~\ref{lem:restricted-line-mixture-bounds} to the conclusion of
	Lemma~\ref{lem:qld-xz-lines} (transported to the sub-line distribution
	via Lemma~\ref{lem:qld-sublines}) and to the consistency relations of
	Lemma~\ref{lem:qld-4-10} --- the resulting quantity is close to
	$\E \sum_{a,b} \bra{\hat\psi}
	\hat{M}^{(\mathrm{Point},X),x}_{a}
	\hat{M}^{(\mathrm{Point},Z),z}_{b} \ot
	\hat{M}^{(\mathrm{Point},Z),z}_{b} \ket{\hat\psi}$,
	which, since $\{\hat{M}^{(\mathrm{Point},X),x}_{a}\}$ is a measurement,
	equals
	$\sum_b \bra{\hat\psi} \hat{M}^{(\mathrm{Point},Z),z}_{b} \ot
	\hat{M}^{(\mathrm{Point},Z),z}_{b} \ket{\hat\psi}
	\ge 1 - \sqrt{\eps}$
	by the self-consistency of the points measurements
	(Lemma~\ref{lem:qld-comm-cons}). The deficit is therefore
	$O(m(\sqrt{\delta_P} + \sqrt{\delta_Q}) + \sqrt{\eps})$, and taking the
	square root gives the claim.
\end{proof}

\begin{proof}[Proof of Lemma~\ref{lem:qld-4-13}]
	\proves{lem:qld-4-13}
	\uses{lem:qld-xz-lines, lem:qld-sublines, def:combine-map,
		lem:qld-4-12, lem:claim-17-1, lem:claim-17-2, lem:claim-17-3}
	Let $D$ be the distribution of Lemma~\ref{lem:qld-sublines} and let
	$D_{|\ell}$ denote $D$ conditioned on the first element being $\ell$.
	For every line $\ell$ in $\F_q^{2m+2}$ and $f \in \deg_{md+1}(\ell)$
	define
	\[
		\hat{Q}^{\ell}_{f}
		= \E_{(\ell_X,\ell_Z) \sim D_{|\ell}}
		\sum_{\substack{f_X \in \deg_{md}(\ell_X),\
			f_Z \in \deg_{md}(\ell_Z):\\
			f = \mathrm{combine}_{f_X,f_Z}}}
		T^{\ell_X,\ell_Z}_{f_X,f_Z}\,,
	\]
	a valid POVM; the polynomial $\mathrm{combine}_{f_X,f_Z}$ of
	Definition~\ref{def:combine-map} is well defined here because Property~2
	of Lemma~\ref{lem:qld-sublines} guarantees that $\ell_X$ and $\ell_Z$
	are compatible with $\ell$. If $\ell$ is axis-parallel, then by
	Property~3 of Lemma~\ref{lem:qld-sublines} so are $\ell_X$ and
	$\ell_Z$, hence $f_X, f_Z \in \deg_d$ by Lemma~\ref{lem:qld-xz-lines};
	along such a line either $(\alpha,\beta)$ is constant, in which case
	$f = \mathrm{combine}_{f_X,f_Z}$ is a fixed linear combination of
	polynomials of degree at most $d$, or $(x,z)$ is constant, in which case
	$f$ is affine in the line parameter, so in either case
	$f \in \deg_d(\ell)$. Chaining Lemmas~\ref{lem:claim-17-1},
	\ref{lem:claim-17-2} and~\ref{lem:claim-17-3} with the triangle
	inequality for scalars gives the consistency bound
	\begin{equation}
		\label{eq:qld-combined-lines-consistency}
		\E_{(\ell,\ell_X,\ell_Z) \sim D}\
		\E_{(x,z,\alpha,\beta) \in \ell}
		\sum_{f_X,f_Z} \bra{\hat\psi}
		\big(T^{\ell_X,\ell_Z}_{f_X,f_Z}\big)_{\mathsf{A}\mathsf{A}'} \ot
		\big(\Id - \hat{Q}^{x,z}_{f_X(x),f_Z(z)}
		\big)_{\mathsf{B}\mathsf{A}''}
		\ket{\hat\psi}
		= O\big(m\,(\eps^{1/4} + \delta_P^{1/4} + \delta_Q^{1/4}
		+ \delta_{\mathrm{Line}}^{1/2})\big)\,;
	\end{equation}
	since
	$\delta_P = \poly(\eps, md/q)$, $\delta_Q = \poly(\eps)$, and
	$\delta_{\mathrm{Line}} = \poly(\eps)$, the right-hand side is a
	function of the form $m \cdot \poly(\eps, md/q)$, and
	$\delta_{\mathrm{combine}}$ is taken to be such a function dominating
	it. (The source states a second bound for the same quantity, whose
	derivation is not supported by the properties proved for the sub-line
	distribution; see Remark~\ref{rem:qld-4-13-source-defects}.) Now
	if the pair $(f_X, f_Z)$ evaluates to $(a, b)$ at $(x, z)$, then
	$\mathrm{combine}_{f_X,f_Z}$ evaluates to $\alpha a + \beta b$ at
	$(x, z, \alpha, \beta)$; hence the consistency of $T$ with
	$\hat{Q}^{x,z}$ along the sub-line pairs expressed by
	\eqref{eq:qld-combined-lines-consistency}, the definition of
	$\hat{Q}^{\ell}$ as the image of $T$ under the combining map, and the
	construction of $\hat{Q}^{x,z,\alpha,\beta}$ in
	Lemma~\ref{lem:qld-4-12} as the corresponding coarse-graining of
	$\hat{Q}^{x,z}$ together yield the first relation of
	\eqref{eq:qld-4-13}. The second relation follows from the analogous
	argument with the registers swapped.
\end{proof}

\begin{remark}[Two error routes in the source proof of
	Lemma~\ref{lem:qld-4-13}]
	\label{rem:qld-4-13-source-defects}
	The source proof in \cite{Ji2020MIPStar} states two different bounds for
	the quantity \eqref{eq:qld-combined-lines-consistency}:
	$O(m(\eps^{1/4} + \delta_P^{1/4} + \delta_Q^{1/4} +
	\delta_{\mathrm{Line}}^{1/2}))$, obtained by chaining the three claims
	(Lemmas~\ref{lem:claim-17-1}--\ref{lem:claim-17-3}), and
	$O(m^2 \delta_P)$, obtained by decomposing the left-hand side of
	\eqref{eq:qld-combined-lines-consistency} into restricted-distribution
	terms and bounding each by \eqref{eq:qld-xz-lines-restricted}; a
	sentence is also duplicated at the corresponding point of the source,
	the second copy with a malformed reference. The second route requires
	the joint law of the pair $((\ell_X, x), (\ell_Z, z))$, for
	$(\ell, \ell_X, \ell_Z) \sim D$ and $(x,z,\alpha,\beta)$ uniform on
	$\ell$, to be a mixture of the product distributions
	$D_{\mathrm{v}_1,i} \times D_{\mathrm{v}_2,j}$, since
	\eqref{eq:qld-xz-lines-restricted} averages over two independently
	sampled line-point pairs; Property~2 of Lemma~\ref{lem:qld-sublines}
	determines only the two marginal laws of $(\ell_X, \ell_Z, x)$ and of
	$(\ell_X, \ell_Z, z)$ separately, not the joint law of
	$((\ell_X, x), (\ell_Z, z))$, so the decomposition invoked by the
	source is not established there. The proof above therefore uses only
	the first route. Neither route, moreover, delivers the error form
	$\poly(m^2\eps, md/q)$ claimed in the source. The first route proves
	the bound $O(m(\eps^{1/4} + \delta_P^{1/4} + \delta_Q^{1/4} +
	\delta_{\mathrm{Line}}^{1/2}))$, a function of the form
	$m \cdot \poly(\eps, md/q)$, and no function of $m^2\eps$ and $md/q$
	alone that vanishes as its arguments vanish dominates it: at
	$\eps = m^{-4}$, and with $q$ taken large, both $m^2\eps$ and $md/q$
	tend to $0$ while the term $m\,\eps^{1/4} = 1$ does not, so the
	derivation does not support the claimed form. Even granting the
	decomposition required by the second route, the bound obtained there
	is $O(m^2\delta_P) = m^2 \cdot \poly(\eps, md/q)$, which again carries
	a separate prefactor polynomial in $m$. The statement of
	Lemma~\ref{lem:qld-4-13} above accordingly records the error as
	$m \cdot \poly(\eps, md/q)$; this form suffices for
	Lemma~\ref{lem:qld-4-7}, whose error function absorbs the prefactor
	$m$ into $(md)^a$. In addition, the final line of
	the source proof of Lemma~\ref{lem:claim-17-2} asserts the bound
	$O(m^2\delta_{\mathrm{Line}})$ for a consistency term that is close to
	one rather than for its deficit; the statements and sketches above
	record the corrected assertions.
\end{remark}

The measurements of Lemma~\ref{lem:qld-4-12} and Lemma~\ref{lem:qld-4-13}
constitute a strategy for the classical low individual degree game in $2m+2$
variables. Applying the quantum soundness of that game yields a single
projective measurement whose outcomes are pairs of global polynomials.

\begin{lemma}[The global polynomial-pair measurement]
	\label{lem:qld-4-7}
	\uses{def:polyfunc, def:expanded-state, def:expanded-point-measurement,
		def:consistency, def:bracket}
	% Hypothesis added: the proof instantiates the classical low individual
	% degree game, and Theorem lem:ld-soundness, at dimension 2m+2, which
	% def:ld-game admits only when the dimension divides q; the source
	% (references/qpbt-paper/14_analysis_of_the_pauli_basis_test.tex,
	% lem:qld-4-7) performs the instantiation without this hypothesis. See
	% rem:qld-4-7-divisibility.
	There exists a function
	$\delta_S(\eps,m,d,q) = a(md)^a(\eps^b + q^{-b} + 2^{-bmd})$, for some
	universal constants $a > 1$ and $0 < b < 1$, such that the following
	holds. Assume that $2m+2$ divides $q$. For
	$\hilb \in \{\hilb_{\mathsf{A}}, \hilb_{\mathsf{B}}\}$ there
	exists a projective measurement $\{\hat{S}_{g_X,g_Z}\}$ acting on
	$\hilb \ot (\C^q)^{\ot n}$, with outcomes consisting of pairs
	$(g_X, g_Z)$ of polynomials, each in $\mathrm{ideg}_{d,m}(\F_q)$, such
	that for $W \in \{X, Z\}$, on average over uniformly random
	$u \in \F_q^m$,
	\begin{align}
		\big(\hat{S}_{[\mathrm{eval}_u(\cdot_W) = a]}
		\big)_{\mathsf{A}\mathsf{A}'}
		&\simeq_{\delta_S}
		\big(\hat{M}^{(\mathrm{Point},W),u}_{a}
		\big)_{\mathsf{B}\mathsf{A}''}\,,
		\label{eq:qld-s-point-con-alice} \\
		\big(\hat{S}_{[\mathrm{eval}_u(\cdot_W) = a]}
		\big)_{\mathsf{B}\mathsf{B}'}
		&\simeq_{\delta_S}
		\big(\hat{M}^{(\mathrm{Point},W),u}_{a}
		\big)_{\mathsf{A}\mathsf{B}''}\,,
		\label{eq:qld-s-point-con-bob}
	\end{align}
	where $\mathrm{eval}_u(\cdot_W)$ denotes the evaluation at the point $u$
	of the first outcome $g_X$ if $W = X$, and of the second outcome $g_Z$
	if $W = Z$.
\end{lemma}
\begin{proof}
	\uses{lem:qld-4-13, lem:qld-4-12, thm:naimark, lem:ld-soundness,
		def:ld-game, fact:agreement, fact:add-a-proj2,
		fact:data-processing, lem:schwartz-zippel-total-degree,
		def:combine-map}
	By Lemmas~\ref{lem:qld-4-12} and~\ref{lem:qld-4-13}, the expanded state
	$\ket{\hat\psi}$ together with the lines measurements
	$\{\hat{Q}^{\ell}_{f}\}$ and the points measurements
	$\{\hat{Q}^{x,z,\alpha,\beta}_{a}\}$ constitutes a strategy for the
	classical low individual degree game (Definition~\ref{def:ld-game}) with
	parameters $(q, 2m+2, d, 1)$ succeeding with probability
	$1 - \delta_{\mathrm{combine}}$; the axis-parallel clause of
	Lemma~\ref{lem:qld-4-13} supplies the degree-$d$ answer format that the
	game requires on axis-parallel lines.
	% The divisibility hypothesis of the statement licenses this
	% instantiation; the source performs it without the hypothesis. See
	% rem:qld-4-7-divisibility.
	(The hypothesis that $2m+2$ divides $q$ makes $(q, 2m+2, d, 1)$ a
	valid parameter tuple for Definition~\ref{def:ld-game}, whose
	dimension parameter is required to divide $q$; admissibility of
	$(q,m,d)$ alone supplies only that $m$ divides $q$, and
	Remark~\ref{rem:qld-4-7-divisibility} records the resulting
	constraint.) Since $\{\hat{Q}^{\ell}_{f}\}$ is
	only a POVM, it is first made projective by Naimark dilation
	(Theorem~\ref{thm:naimark}), which preserves the $\simeq$ conclusions of
	Lemma~\ref{lem:qld-4-13}. The quantum soundness of the classical test
	(Theorem~\ref{lem:ld-soundness}, applied with $m$ replaced by $2m+2$ and
	$\eps$ by $\delta_{\mathrm{combine}}$) then yields a POVM
	$\{\hat{S}_g\}$ with outcomes $g \in \mathrm{ideg}_{d,2m+2}(\F_q)$ that
	is $\delta_{\mathrm{ld}}$-self-consistent and
	$\delta_{\mathrm{ld}}$-consistent with the points measurements
	$\{\hat{Q}^{x,z,\alpha,\beta}_{a}\}$, where
	$\delta_{\mathrm{ld}}
	= a(md)^a(\delta_{\mathrm{combine}}^b + q^{-b} + 2^{-bmd})$
	for universal constants $a, b$; a second application of Naimark dilation
	makes $\{\hat{S}_g\}$ projective with the same guarantees. Let
	$\mathcal{G}$ be the set of polynomials of the form
	$g(x,z,\alpha,\beta) = \alpha g_X(x) + \beta g_Z(z)$ with
	$g_X, g_Z \in \mathrm{ideg}_{d,m}(\F_q)$. A chain of approximations
	built from the consistency of $\{\hat{S}_g\}$ with the points
	measurements and the relations
	\eqref{eq:qld-4-12-cons-m-hat}--\eqref{eq:qld-4-12-cons-m-hat2} of
	Lemma~\ref{lem:qld-4-12} (using Lemma~\ref{fact:agreement} to convert
	$\simeq$ into $\approx$ and Lemma~\ref{fact:add-a-proj2}) correlates
	$\hat{S}_g$ with the indicator-weighted products of the $X$ and $Z$
	point measurements in both orders. Combining this with Schwartz--Zippel
	case analyses (Lemma~\ref{lem:schwartz-zippel-total-degree}, applied to
	polynomials of total degree at most $(2m+2)d$) --- first ruling out the
	$g$ that are not linear in $(\alpha, \beta)$, then, among the linear
	ones $g = \alpha g_1(x,z) + \beta g_2(x,z)$, ruling out those whose
	coefficient polynomials $g_1$ or $g_2$ depend on the wrong block of
	variables --- yields
	$\sum_{g \notin \mathcal{G}}
	\| (\hat{S}_g)_{\mathsf{A}\mathsf{A}'} \ket{\hat\psi} \|^2
	\le O(\delta_{\mathrm{ld}} + \delta_Q + md/q)
	=: \delta_{\mathcal{G}}$.
	Define the projective sub-measurement
	$\hat{S}_{g_X,g_Z} = \hat{S}_{\mathrm{combine}_{g_X,g_Z}}$
	(Definition~\ref{def:combine-map}); its incompleteness is at most
	$\delta_{\mathcal{G}}$, and a further approximation chain with
	successive errors $\delta_{\mathcal{G}}$, $\delta_Q^{1/2}$,
	$\delta_Q^{1/2}$, $1/q$, $1/q$ gives
	\[
		\big(\hat{S}_{[\mathrm{eval}_x(\cdot) = a,\
			\mathrm{eval}_z(\cdot) = b]}
		\big)_{\mathsf{A}\mathsf{A}'}
		\simeq_{O(\delta_{\mathcal{G}} + \delta_Q^{1/2} + 1/q)}
		\big(\hat{M}^{(\mathrm{Point},X),x}_{a}\,
		\hat{M}^{(\mathrm{Point},Z),z}_{b}\,
		\hat{M}^{(\mathrm{Point},X),x}_{a}
		\big)_{\mathsf{B}\mathsf{A}''}\,.
	\]
	Completing the sub-measurement by adding
	$\Id - \sum_{g_X,g_Z} \hat{S}_{g_X,g_Z}$ to an arbitrary element costs a
	further $\delta_{\mathcal{G}}$, and the resulting error
	$\delta_S = O(\delta_{\mathrm{ld}} + \delta_Q^{1/2} + md/q)$ has the
	stated form for suitable universal constants (the prefactor $m$ in
	$\delta_{\mathrm{combine}} = m \cdot \poly(\eps, md/q)$ is absorbed
	into $(md)^a$ after enlarging $a$ and shrinking $b$); an application of
	Lemma~\ref{fact:data-processing} then gives
	\eqref{eq:qld-s-point-con-alice} for $W = X$. For $W = Z$ the same
	steps are performed with $X$ and $Z$ interchanged, starting from the
	opposite-order correlation, and entirely analogous calculations with
	the registers swapped yield \eqref{eq:qld-s-point-con-bob}.
\end{proof}

% This remark records a defect of the source argument; it is not a statement
% from the paper.
\begin{remark}[Divisibility hypothesis in the application of the classical
	soundness theorem]
	\label{rem:qld-4-7-divisibility}
	Definition~\ref{def:ld-game} defines the classical low individual
	degree game only when its dimension parameter divides $q$: the index
	function $\chi$ of the question distribution
	(Definition~\ref{def:ld-question-distribution}) partitions $\F_q$ into
	classes of equal size, one for each coordinate. The same divisibility
	underlies the line-point distribution over $\F_q^{2m+2}$
	(Definition~\ref{def:line-point-dist} instantiated at dimension
	$2m+2$). For this reason Lemmas~\ref{lem:qld-sublines},
	\ref{lem:qld-4-13} and~\ref{lem:qld-4-7} carry the explicit hypothesis
	that $2m+2$ divides $q$, under which their proofs go through as
	written. This hypothesis does not follow from admissibility of
	$(q,m,d)$ (Definition~\ref{def:admissible}), which supplies only that
	$m$ divides $q$; indeed the two conditions are jointly satisfiable
	only in a degenerate range: an admissible field size
	(Definition~\ref{def:admissible-size}) is a power of two, so $m$
	dividing $q$ forces $m$ to be a power of two, while
	$2m+2 = 2(m+1)$ divides the power of two $q$ only if $m+1$ is also a
	power of two, and both conditions hold only for $m = 1$ (in which case
	$2m+2 = 4$ divides every admissible $q \ge 4$). The source proof in
	\cite{Ji2020MIPStar} makes the instantiation at dimension $2m+2$
	without addressing the divisibility. Consequently, for admissible
	parameters with $m \ge 2$ the hypothesis fails, the statement of
	Lemma~\ref{lem:qld-4-7} is vacuous, and the conclusion of that lemma
	is not established for such parameters, here or in the source; a
	proof covering every admissible parameter tuple requires a
	formulation of the game, of its question distribution, and of
	Theorem~\ref{lem:ld-soundness} that applies to dimensions not dividing
	$q$.
\end{remark}

\begin{remark}[Constant regimes in the quotation of the classical soundness
	theorem]
	\label{rem:qld-4-7-constants}
	Theorem~\ref{lem:ld-soundness} states its error function with universal
	constants $a \ge 1$ and $0 < b \le 1$, while the source proof of
	Lemma~\ref{lem:qld-4-7} in \cite{Ji2020MIPStar} quotes it, and the
	statement of Lemma~\ref{lem:qld-4-7} is phrased, with $a > 1$ and
	$0 < b < 1$. The discrepancy is harmless: enlarging $a$ and shrinking
	$b$ only weakens a bound of this form, so any admissible pair for
	Theorem~\ref{lem:ld-soundness} yields an admissible pair for
	Lemma~\ref{lem:qld-4-7}. The closing computation of the source proof
	also mixes the terms $d/q$ and $md/q$ in the bound for $\delta_S$; the
	weaker $md/q$ is the one recorded above.
\end{remark}
